
Or
prenons pour @@M0@@ une \uline{variété de drapeaux}
@@M1@@
où @@M2@@ est un groupe algébrique semi-simple simplement connexe sur le corps algébriquement clos @@M3@@, et @@M4@@ un sous-groupe de Borel [5]. Lorsque @@M5@@, alors la décomposition cellulaire de Bruhat de @@M6@@ [5, Exp. 13] implique que l'homomorphisme canonique
@@M7@@
est un isomorphisme, de même d'ailleurs que l'homomorphisme canonique
@@M8@@
où @@M9@@ est l'espace compact des points complexes de @@M10@@, muni de sa topologie habituelle déduite de la topologie de @@M11@@. D'autre part, il a été prouvé par L. HODGKIN [12] que @@M12@@ est bien engendré par les classes des faisceaux inversibles sur @@M13@@. Il en est donc de même de @@M14@@ lui-même, et il n'est d'ailleurs pas difficile d'en déduire, par une technique de spécialisation utilisant les groupes semi-simples déployés sur @@M15@@ (SGA 3), que le même résultat reste valable pour tout corps @@M16@@ algébriquement clos. D'ailleurs, la décomposition cellulaire de @@M17@@ (ou la bijectivité de @@M18@@ dans le cas @@M19@@) implique aisément que @@M20@@ est sans torsion, donc grâce à 4.2 l'homomorphisme @@M21@@ est un isomorphisme. De ces faits il résulte que @@M22@@ est un isomorphisme si et seulement si @@M23@@ est engendré par ses éléments de degré @@M24@@, ou encore, si @@M25@@, si et seulement si la cohomologie entière de @@M26@@ est engendrée par les éléments de degré @@M27@@, i.e. si le classifiant @@M28@@ est sans torsion. Il est connu d'autre part [15, p. 217 th. B (1)@@M29@@(2), et 2.5] que ceci est le cas si et seulement si @@M30@@ est "sans torsion", ou encore si @@M31@@ est isomorphe à un produit de groupes @@M32@@ et @@M33@@. Il suffit donc de prendre pour @@M34@@ un groupe @@M35@@ avec @@M36@@, ou un groupe simple exceptionnel, pour avoir un exemple où @@M37@@ n'est pas un isomorphisme.

\textbf{4.6.} On notera que dans cet exemple, il est hors de doute que le "bon" anneau est l'anneau sans torsion @@M38@@, ayant comme base les cellules de la décomposition de Bruhat, anneau également isomorphe à la cohomologie entière de @@M39@@ si @@M40@@; par contre, l'anneau @@M41@@ apparaît comme un anneau ayant de la torsion de nature assez pathologique. Cet exemple confirme donc la conception suivant laquelle l'invariant @@M42@@ n'est guère
